\documentclass[11pt,a4paper]{article}
\usepackage[T1]{fontenc}
\usepackage[utf8]{inputenc}
\usepackage{amsmath,amssymb,amsthm}
\usepackage[margin=1in]{geometry}
\usepackage[colorlinks=true,linkcolor=blue,urlcolor=blue]{hyperref}
\theoremstyle{plain}
\newtheorem{theorem}{Theorem}
\newtheorem{lemma}[theorem]{Lemma}
\newtheorem{proposition}[theorem]{Proposition}
\newtheorem{corollary}[theorem]{Corollary}
\theoremstyle{definition}
\newtheorem{remark}[theorem]{Remark}

\title{The empirical recurrence for OEIS A208026, proved by transfer matrix}
\author{Adrian Perez Fontelles\\ \small Independent researcher}
\date{31 August 2026}
\begin{document}
\maketitle

\begin{abstract}
OEIS A208026 counts arrays over $\{0,\dots,1\}$ of width $6$ that avoid the patterns
$(0,0,0), (0,0,1)$ along one axis and $(0,0,0), (1,0,1)$ along the other, and carries an empirical recurrence of
order $81$ contributed by R.~H.~Hardin. It is true. The patterns have length $3$, so
the constraint across the growing direction ties together $3$ consecutive rows: taking a
window of $2$ consecutive rows as a state makes the array a walk in a finite digraph
on $S=676$ vertices, and $a(n)=\mathbf{1}^{\!\top}P^{\,n-2}\mathbf{1}$. Writing $q$
for the characteristic polynomial of the conjectured recurrence, the residual is
$\mathbf{1}^{\!\top}P^{\,n-2-81}q(P)\mathbf{1}$, which vanishes identically
beyond a threshold that the computation pins down exactly. Finite, exact, integer.
\end{abstract}

\noindent\small 2020 Mathematics Subject Classification. 05A15, 68R15, 15B36.\normalsize

\section{The sequence and the conjecture}

OEIS A208026 is ``Number of nX6 0..1 arrays avoiding 0 0 0 and 0 0 1 horizontally and 0 0 0 and 1 0 1 vertically.''. It has offset
$1$ and begins
\[
26, 676, 4374, 26685, 158847, 955073, 5788524,\ \dots
\]
The entry carries the following comment:
\begin{quote}\small
Empirical: a(n) = a(n-1) +3*a(n-2) +106*a(n-3) +305*a(n-4) +688*a(n-5) -487*a(n-6) -7636*a(n-7) -22388*a(n-8) -16653*a(n-9) +105705*a(n-10) +368317*a(n-11) +296151*a(n-12) -1322186*a(n-13) -3797589*a(n-14) -461119*a(n-15) +12216056*a(n-16) +16314675*a(n-17) -11652578*a(n-18) -48384768*a(n-19) -31494807*a(n-20) +52484033*a(n-21) +108587927*a(n-22) +37697973*a(n-23) -108118565*a(n-24) -174749969*a(n-25) -53159946*a(n-26) +135111621*a(n-27) +237174429*a(n-28) +96181263*a(n-29) -81439514*a(n-30) -277585453*a(n-31) -182415247*a(n-32) +25256104*a(n-33) +180399360*a(n-34) +347994520*a(n-35) -90243930*a(n-36) +165953626*a(n-37) -662537687*a(n-38) +556623889*a(n-39) -991152494*a(n-40) +1372650960*a(n-41) -1621551426*a(n-42) +1909898571*a(n-43) -2182104283*a(n-44) +2186815242*a(n-45) -2123785604*a(n-46) +2065891574*a(n-47) -1695531591*a(n-48) +1513432468*a(n-49) -1178637615*a(n-50) +918618751*a(n-51) -671511358*a(n-52) +497302975*a(n-53) -324440300*a(n-54) +223677440*a(n-55) -145261693*a(n-56) +84442881*a(n-57) -54971856*a(n-58) +30801717*a(n-59) -17151395*a(n-60) +9895586*a(n-61) -5453123*a(n-62) +2470341*a(n-63) -1537394*a(n-64) +757067*a(n-65) -221697*a(n-66) +233701*a(n-67) -62539*a(n-68) +16901*a(n-69) -29422*a(n-70) +1837*a(n-71) -2160*a(n-72) +2083*a(n-73) +61*a(n-74) +274*a(n-75) -44*a(n-76) +10*a(n-77) -24*a(n-78) -4*a(n-79) +a(n-81) for n>82
\end{quote}
As of the ``Last modified'' line on the live entry (Oct 03 2025 14:34:55, revision
9) the statement is still recorded as empirical, and nothing on the entry
records it as settled either way.

Written out, the claim is
\begin{equation}\label{eq:conj}
a(n)\;=\;a(n-1) + 3\,a(n-2) + 106\,a(n-3) + 305\,a(n-4) + 688\,a(n-5) - 487\,a(n-6) - 7636\,a(n-7) - 22388\,a(n-8) - 16653\,a(n-9) + 105705\,a(n-10) + 368317\,a(n-11) + 296151\,a(n-12) - 1322186\,a(n-13) - 3797589\,a(n-14) - 461119\,a(n-15) + 12216056\,a(n-16) + 16314675\,a(n-17) - 11652578\,a(n-18) - 48384768\,a(n-19) - 31494807\,a(n-20) + 52484033\,a(n-21) + 108587927\,a(n-22) + 37697973\,a(n-23) - 108118565\,a(n-24) - 174749969\,a(n-25) - 53159946\,a(n-26) + 135111621\,a(n-27) + 237174429\,a(n-28) + 96181263\,a(n-29) - 81439514\,a(n-30) - 277585453\,a(n-31) - 182415247\,a(n-32) + 25256104\,a(n-33) + 180399360\,a(n-34) + 347994520\,a(n-35) - 90243930\,a(n-36) + 165953626\,a(n-37) - 662537687\,a(n-38) + 556623889\,a(n-39) - 991152494\,a(n-40) + 1372650960\,a(n-41) - 1621551426\,a(n-42) + 1909898571\,a(n-43) - 2182104283\,a(n-44) + 2186815242\,a(n-45) - 2123785604\,a(n-46) + 2065891574\,a(n-47) - 1695531591\,a(n-48) + 1513432468\,a(n-49) - 1178637615\,a(n-50) + 918618751\,a(n-51) - 671511358\,a(n-52) + 497302975\,a(n-53) - 324440300\,a(n-54) + 223677440\,a(n-55) - 145261693\,a(n-56) + 84442881\,a(n-57) - 54971856\,a(n-58) + 30801717\,a(n-59) - 17151395\,a(n-60) + 9895586\,a(n-61) - 5453123\,a(n-62) + 2470341\,a(n-63) - 1537394\,a(n-64) + 757067\,a(n-65) - 221697\,a(n-66) + 233701\,a(n-67) - 62539\,a(n-68) + 16901\,a(n-69) - 29422\,a(n-70) + 1837\,a(n-71) - 2160\,a(n-72) + 2083\,a(n-73) + 61\,a(n-74) + 274\,a(n-75) - 44\,a(n-76) + 10\,a(n-77) - 24\,a(n-78) - 4\,a(n-79) + a(n-81)
\end{equation}
for all $n>81$. 

\section{The array as a walk}

Call a row \emph{admissible} if the word it forms contains none of the patterns
$(0,0,0), (0,0,1)$ as $3$ consecutive entries. There are $26$ admissible rows; that condition
involves a single row only.

The remaining patterns $(0,0,0), (1,0,1)$ are forbidden along the other axis, so each of them
constrains $3$ consecutive rows at a fixed position of the column. Take as state a
window $w=\bigl(w^{(1)},\dots,w^{(2)}\bigr)$ of $2$ consecutive admissible
rows, so $S=676$, and declare $w\to w'$ an edge when $w'$ is $w$ with its first entry
dropped and one further admissible row $t$ appended, subject to: for every position
$j$, the $3$-tuple read down the column $j$ across $w^{(1)},\dots,w^{(2)},t$
is none of $(0,0,0), (1,0,1)$. Let $P$ be the adjacency matrix.

\begin{lemma}\label{lem:walk}
$a(n)=\mathbf{1}^{\!\top}P^{\,n-2}\mathbf{1}$ for every $n\ge2$.
\end{lemma}

\begin{proof}
An array of the entry's kind with $n$ rows is a sequence of $n$ admissible rows in
which no forbidden pattern appears along the other axis. Reading it as the sequence of its
overlapping windows of length $2$ turns it into a walk: consecutive windows overlap in
$1$ rows by construction, and the edge condition is exactly the requirement on the one
new $3$-tuple that each step creates. Conversely a walk reassembles into such an array,
and the two constructions are mutually inverse. A walk of length $n-2$ visits $n-1$
windows, that is $n$ rows; summing over all starting and ending states gives
$\mathbf{1}^{\!\top}P^{\,n-2}\mathbf{1}$.
\end{proof}

\section{The criterion}

Let $q(t)=t^{81}-\sum_i c_i t^{81-i}$ be the characteristic polynomial of
\eqref{eq:conj} and $u_j=\mathbf{1}^{\!\top}P^{j}q(P)\mathbf{1}$.

\begin{lemma}\label{lem:crit}
$a(n)-\sum_i c_i a(n-i)=u_{\,n-2-81}$ whenever $n-81\ge2$. Hence
\eqref{eq:conj} holds for every $n>t$ exactly when $u_j=0$ for every
$j>t-2-81$. Moreover $(u_j)$ satisfies the monic linear recurrence whose
characteristic polynomial is that of $P$, of degree $S$, so $S$ consecutive zeros force all
later ones and the last nonzero $u_j$ pins $t$ down exactly.
\end{lemma}

\begin{proof}
Substituting Lemma~\ref{lem:walk} into the left side and factoring out
$P^{\,n-2-81}$ gives the identity. The equivalence follows by letting the exponent
run. For the last claim, Cayley--Hamilton gives $\chi(P)=0$ for the characteristic
polynomial $\chi$ of $P$, so $\chi$ applied to the shift annihilates $(u_j)$; being monic of
degree $S$, it propagates $S$ consecutive zeros forward.
\end{proof}

\section{The computation}

\begin{theorem}
The recurrence \eqref{eq:conj} holds for every $n>81$. This is the statement of
Section~1.
\end{theorem}

\begin{proof}
The vector $q(P)\mathbf{1}\in\mathbb{Z}^{S}$ was formed by $81$ matrix--vector
products in exact integer arithmetic and the $u_j$ evaluated in turn. They vanish from the
index corresponding to $n=81+1$ onwards, and the run was continued past $S$ consecutive
zeros, which by Lemma~\ref{lem:crit} settles every larger index.
\end{proof}

\section{Verification}

Three checks, each independent of the algebra above.

First, the digraph was built directly from the entry's stated patterns and the walk counts
compared against the entry's DATA at every one of the $19$ published indices; the
values agree exactly, including the initial count of $26$ admissible rows. That is
what ties the digraph to the sequence, and a model reproducing only part of the published
data was rejected wherever it occurred.

Second, the recurrence \eqref{eq:conj} was evaluated on those published terms in exact
integer arithmetic, with no matrices involved, and vanishes at every index the entry claims.

Third, the vanishing test of Lemma~\ref{lem:crit} was carried past $S=676$ consecutive
zeros rather than to a sampled prefix.

\begin{thebibliography}{9}
\bibitem{oeis} The OEIS Foundation, \emph{The On-Line Encyclopedia of Integer
Sequences}, \url{https://oeis.org}, 2026. Sequence A208026.
\bibitem{stanley} R.~P.~Stanley, \emph{Enumerative Combinatorics}, Volume 1, 2nd ed.,
Cambridge University Press, 2011. (Transfer-matrix method, Section 4.7.)
\bibitem{fs} P.~Flajolet and R.~Sedgewick, \emph{Analytic Combinatorics}, Cambridge
University Press, 2009.
\bibitem{horn} R.~A.~Horn and C.~R.~Johnson, \emph{Matrix Analysis}, 2nd ed., Cambridge
University Press, 2013.
\end{thebibliography}

\end{document}
